% Part:sets-functions-relations
% Chapter: size-of-sets
% Section: non-enumerability-alt

\documentclass[../../../include/open-logic-section]{subfiles}

\begin{document}

\olfileid{sfr}{siz}{nen-alt}

\olsection{\printtoken{S}{nonenumerable} سټونه}

\begin{editorial}
  دا برخه په \olref[enm-alt]{sec} کښې د تعريفونو په کارولو، يعنې
  له $\PosInt$ څخه د پر هدف سټ بشپړې تابع پر ځاے له~$\Nat$ سره
  د دوه اړخيزې يو پر يو تابع په غوښتلو، د $\Bin^\omega$ او
  $\Pow{\Nat}$ د شمېر نۀ وړ والے ثابتوي۔
\end{editorial}

\begin{explain}
د طبيعي عددونو سټ $\Nat$ نامتناهي دے۔ دا په څرګند ډول
!!{enumerable} هم دے۔ خو د پام وړ حقيقت دا دے چې
\emph{!!{nonenumerable}} سټونه، يعنې هغه سټونه چې !!{enumerable}
نۀ دي، شته (\olref[sfr][siz][enm-alt]{defn:enumerable} وګورئ)۔

دا ښايي حيرانوونکې وي۔ بالاخره، دا ويل چې $A$ !!{nonenumerable}
دے، دا ويل دي چې \emph{هېڅ} !!{bijection}
$f \colon \Nat \to A$ نۀ شته؛ يعنې هېڅ هغه تابع چې د~$\Nat$
بې‌شمېره !!{element}s پر~$A$ نقشوي، ټول~$A$ نۀ شي پوره کولے۔
نو که $A$ !!{nonenumerable} وي، د طبيعي عددونو په پرتله د~$A$
«زيات» !!{element}s شته۔

د دې د ثابتولو دپاره چې يو سټ !!{nonenumerable} دے، بايد وښيئ چې هېڅ
مناسبه !!{bijection} نۀ شي موجودېدلای۔ د دې تر ټولو ښه لار دا ده
چې وښيئ د~$A$ د !!{element}s د شمېرلو هره هڅه بايد لږ تر لږه يو
!!{element} پرېږدي؛ دا ښيي چې هېڅ تابع $f\colon \Nat \to A$
!!{surjective} نۀ ده۔ د دې د ثابتولو يوه عمومي تګلاره د کانتور
\emph{قطري طريقه} ده۔ د~$A$ د !!{element}s يو لړ، د بېلګې په
ډول $x_1$، $x_2$، \dots، چې راکړل شي، موږ د~$A$ يو بل
!!{element} جوړوو چې د خپل جوړښت له مخې په هېڅ ډول په هغه لړ
کښې نۀ شي راتلے۔

خو دا هر څه په بېلګه تر ټولو ښه پوهېدل کېږي۔ نو زموږ لومړۍ
بېلګه د $0$ګانو او $1$ګانو د ټولو نامتناهي توريزو تسلسلونو سټ
~$\Bin^\omega$ دے۔ (`$\Bin$' د دوه‌ګوني دپاره دے، او يوازې ئې د
دوه غړيز سټ $\{0,1\}$ په توګه ګڼلے شو۔)\oliflabeldef{sfr:card-arithmetic:card-opps:sec}{\footnote{په لا دقيقه توګه، بايد وټاکو چې
$\Bin^\omega$ د $0$ګانو او $1$ګانو د ټولو $\omega$-تسلسلونو
سټ، يعنې سټ $\funfromto{\omega}{\{0,1\}}$، دے۔ خو د دې معنا به
يوازې په \olref[sfr][card-arithmetic][card-opps]{sec} کښې روښانه
شي۔}{} خو اوس مهال دا لږ نرم بيان بايد هېڅ ګډوډي جوړه نۀ کړي۔}
\end{explain}

\begin{thm}
\ollabel{thm:nonenum-bin-omega}
$\Bin^\omega$~!!{nonenumerable} دے۔
\end{thm}

\begin{proof}
د $\Bin^\omega$ د يوه فرعي سټ هره شمېرنه په پام کښې ونيسئ۔ نو
کوم لړ $s_{0}$، $s_{1}$، $s_{2}$، \dots{} لرو چې هر $s_n$ د
$0$ګانو او~$1$ګانو يو نامتناهي توريز تسلسل دے۔ فرض کړئ
$s_n(m)$ په دې لړ کښې د $n$م توريز تسلسل $m$م رقم دے۔
\textbf{د ژباړې د نسخې سمون (OLSIZ-008):} انګرېزي نثر د تسلسل
او رقم د شاخصونو رولونه سرچپه بيان کړي وو؛ لاندې اصلي جدول لومړے
شاخص د قطار او دويم شاخص د ستون په توګه ټاکي۔ نو اوس خپل لړ د يوه جدولي
ترتيب په توګه ګڼلے شو، چې $s_n(m)$ په $n$م قطار او $m$م ستون
کښې اېښوول شوے دے:
\[
\begin{array}{c|c|c|c|c|c}
& 0 & 1 & 2 & 3 & \dots \\\hline
0 & \mathbf{s_{0}(0)} & s_{0}(1) & s_{0}(2) & s_0(3) & \dots \\\hline
1 & s_{1}(0)& \mathbf{s_{1}(1)} & s_1(2) & s_1(3) & \dots \\\hline
2 & s_{2}(0)& s_{2}(1) & \mathbf{s_2(2)} & s_2(3) & \dots \\\hline
3 & s_{3}(0)& s_{3}(1) & s_3(2) & \mathbf{s_3(3)} & \dots \\\hline
\vdots & \vdots & \vdots & \vdots & \vdots & \mathbf{\ddots}
\end{array}
\]
اوس د $0$ګانو او $1$ګانو يو داسې نامتناهي توريز تسلسل $d$
جوړوو چې په دې لړ کښې نۀ وي۔ دا د هغه د هرې خانې په ټاکلو؛
يعنې د ټولو $n \in \Nat$ دپاره د $d(n)$ په ټاکلو، کوو۔ په شهودي
ډول، د پورته جدول پر قطر لاندې ځو (ځکه ورته «قطري طريقه» وايو)،
او بيا هر $1$ په $0$ او هر $0$ په~$1$ بدلوو۔
\textbf{د ژباړې د نسخې سمون (OLSIZ-009):} انګرېزي نثر د دويم
بدلون پر ځاے لومړے بدلون بيا تکرار کړے و؛ لاندې اصلي حالت‌تعريف
د صفر او يو دواړه اړخونه ټاکي۔ په لا مجرد ډول، $d(n)$ د دې له
مخې $0$ يا $1$ تعريفوو چې د قطر $n$م !!{element}، $s_n(n)$،
$1$ دے که $0$؛ يعنې:
\[
d(n) =
\begin{cases}
1 & \text{که $s_{n}(n) = 0$ وي}\\
0 & \text{که $s_{n}(n) = 1$ وي}
\end{cases}
\]
په څرګند ډول $d \in \Bin^\omega$، ځکه دا د $0$ګانو او $1$ګانو
يو نامتناهي توريز تسلسل دے۔ خو موږ $d$ داسې جوړ کړے چې د هر
$n \in \Nat$ دپاره $d(n) \neq s_n(n)$۔ يعنې $d$ له $s_n$ څخه
په خپل $n$م مقام کښې توپير لري۔ نو د هر $n\in \Nat$ دپاره
$d \neq s_n$۔ ځکه $d$ په لړ $s_0$، $s_1$، $s_2$،
\dots کښې نۀ شي راتلے۔

موږ وښودل چې د $\Bin^\omega$ د کوم فرعي سټ هره خپله سرې شمېرنه
به د $\Bin^\omega$ کوم !!{element} پرېږدي۔ نو د سټ
$\Bin^\omega$ هېڅ شمېرنه نۀ شته؛ يعنې $\Bin^\omega$
!!{nonenumerable} دے۔
\end{proof}

\begin{explain}
د ثبوت دې طريقې ته «قطرول» ويل کېږي، ځکه د~$d$ د تعريف دپاره
د جدول قطر کاروي۔ خو قطرول د جدول شتون ته اړتيا نۀ لري۔ په
حقيقت کښې، په ورته مفکوره ښودلے شو چې ځينې سټونه
!!{nonenumerable} دي، ان چې هېڅ جدول او رښتينے قطر نۀ وي۔
لاندې پايله ښيي چې څنګه۔
\end{explain}

\begin{thm}
\ollabel{thm:nonenum-pownat}
$\Pow{\Nat}$ !!{enumerable} نۀ دے۔
\end{thm}

\begin{proof}
په هماغه طريقه مخکښې ځو: ښيو چې د~$\Nat$ د فرعي سټونو هر لړ
د $\Nat$ کوم فرعي سټ پرېږدي۔ نو فرض کړئ چې د $\Nat$ د فرعي
سټونو کوم لړ $N_0, N_1, N_2, \ldots$ لرو۔ سټ $D$ داسې تعريفوو:
$n \in D$ هله او يوازې هله چې $n \notin N_{n}$:
\[
D = \Setabs{n \in \Nat}{n \notin N_n}
\]
په څرګند ډول $D\subseteq \Nat$۔ خو $D$ په لړ کښې نۀ شي راتلے۔
بالاخره، د جوړښت له مخې $n \in D$ هله او يوازې هله چې
$n\notin N_n$؛ نو د هر $n \in \Nat$ دپاره $D \neq N_n$۔
\end{proof}

\begin{explain}
تېر ثبوت د قطر نوم وانۀ خيست۔ بيا هم که داسې ئې انځور کړئ، قطري
طريقه په کښې ليدلے شئ: وانګېړئ چې سټونه $N_0$، $N_1$، \dots په
جدولي ترتيب کښې ليکل شوي، چې $N_n$ په $n$م قطار کښې داسې ليکو
چې $m$ په $m$م ستون کښې هله او يوازې هله ليکو که $m \in N_n$۔
د بېلګې په ډول، فرض کړئ په هغه لړ کښې لومړي څلور سټونه
$\{0,1,2,\dots\}$، $\{1, 3, 5, \dots\}$، $\{0,1,4\}$ او
$\{2,3,4,\dots\}$ دي؛ نو زموږ جدولي ترتيب به داسې پيل شي:
\[
\begin{array}{r@{}rrrrrrr}
  N_0 = \{ & \mathbf{0}, & 1, & 2, & & & & \dots\}\\
  N_1 = \{ &  & \mathbf{1}, &  & 3, &  & 5, & \dots\}\\
  N_2 = \{ & 0, & 1, &  &  & 4\phantom{,} &  & \}\\
  N_3 = \{ &  &  & 2, & \mathbf{3}, & 4, & & \dots\}\\
  &\vdots & & & & & & \ddots\phantom{\}}
  \end{array}
\]
بيا $D$ هغه سټ دے چې پر قطر د لاندې تلو له لارې تر لاسه کېږي:
$n \in D$ هله او يوازې هله اېږدو چې $n$ پر قطر \emph{نۀ} وي۔
نو په پورته حالت کښې $0$ او $1$ پرېږدو،~$2$ شاملوو،~$3$
پرېږدو، او همداسې نور۔
\end{explain}

\begin{prob}
په څرګند قطري استدلال وښيئ چې د ټولو تابعو
$f \colon \Nat \to \Nat$ سټ !!{nonenumerable} دے۔ يعنې وښيئ
چې که $f_1$، $f_2$، \dots د تابعو يو لړ وي او هره
$f_i\colon \Nat \to \Nat$ وي، نو کومه $g \colon \Nat \to
\Nat$ شته چې په دې لړ کښې نۀ ده۔
\end{prob}

\end{document}
